
\documentclass[a4paper,10pt]{report}

%\usepackage[landscape]{geometry}
\usepackage[cm-default]{fontspec}
\setromanfont{FreeSerif}
\setsansfont{FreeSans}
\setmonofont{FreeMono}
\usepackage[utf8x]{inputenc}
\usepackage{fontspec}
\usepackage{xunicode}
\usepackage{xltxtra}
\usepackage{xgreek}
\usepackage{amsmath}
\usepackage{unicode-math}
\usepackage{ulem}
\usepackage{color}
\usepackage{verbatim}
\usepackage{nopageno}
\usepackage{graphicx}
\usepackage{textpos}
\setlength{\TPHorizModule}{1cm}
\setlength{\TPVertModule}{1cm}
%\usepackage[colorgrid,texcoord]{eso-pic}
\usepackage[outline]{contour}
\usepackage{wrapfig}
\usepackage{url}
\usepackage{color}
\usepackage{tikz}
\usepackage{fancybox,fancyhdr}
\usepackage{subfigure}
\usepackage{pstricks}
\usepackage{epsfig}
\usepackage{multicol}
\usepackage{listings}
\usepackage{enumerate}
\usepackage{hyperref}
\hypersetup{
  bookmarks=true,
  bookmarksopen=true,
  pdfborder=false,
  pdfpagemode=UseNone,
  raiselinks=true,
  pdfhighlight={/P},
  colorlinks,
  citecolor=black,
  filecolor=black,
  linkcolor=black,
  urlcolor=black
}


\usepackage{fontspec}
\usepackage{xunicode}
\usepackage{xltxtra}


% Margins
%\setlength{\textwidth}{24cm}

%\setlength{\voffset}{0in}
%\setlength{\textheight}{6.8in}
% Colors
\definecolor{cbrown}{rgb}{0.49,0.24,0.07}
\definecolor{cmpez}{rgb}{0.92,0.88,0.79}
\definecolor{cmpez2}{rgb}{0.62,0.33,0.34}
\definecolor{cwhite}{rgb}{1,1,1}
\definecolor{cblack}{rgb}{0,0,0}
\definecolor{cred}{rgb}{0.9,0.15,0.15}
\definecolor{clblue}{rgb}{0.57,0.85,0.97}
\definecolor{corange}{rgb}{0.86,0.49,0.18}
\definecolor{clorange}{rgb}{0.95,0.84,0.61}
\definecolor{cyellow}{rgb}{1,0.95,0.16}
\definecolor{cgreen}{rgb}{0.69,0.8,0.31}
\definecolor{clbrown}{rgb}{0.78,0.67,0.43}
\definecolor{cmagenta}{rgb}{0.79,0.05,0.54}
\definecolor{cgray}{rgb}{0.76,0.73,0.63}
\definecolor{bluesite}{rgb}{0.05,0.43,0.67}
\definecolor{redsite}{rgb}{0.90,0.15,0.15}
\definecolor{backsite}{rgb}{0.07,0.07,0.07}
% ----------- from Costas -----------------------------

%---------------   Main Fonts    -----------------%   

% \setmainfont[Mapping=tex-text]{Calibri}
% \setmainfont[Mapping=tex-text]{Times New Roman}
%\setmainfont[Mapping=tex-text]{Droid Serif}
%\setmainfont[Mapping=tex-text]{Cambria}
%\setmainfont[Mapping=tex-text]{cm-unicode}
% \setmainfont[Mapping=tex-text]{Gentium}
% \setmainfont[Mapping=tex-text]{GFS Didot}
% \setmainfont[Mapping=tex-text]{Comic Sans MS}
 \setmainfont[Mapping=tex-text]{Ubuntu}
% \setmainfont[Mapping=tex-text]{Myriad Pro}
% \setmainfont[Mapping=tex-text]{CMU Concrete}
%\setmainfont[Mapping=tex-text]{DejaVu Sans}
%\setmainfont[Mapping=tex-text]{KerkisSans}
%\setmainfont[Mapping=tex-text]{KerkisCaligraphic}
% \setmainfont[Mapping=tex-text]{Segoe Print}
% \setmainfont[Mapping=tex-text]{Gabriola}

%------------------------------------------------%

% ------------ Mathematics Fonts ----------------%
\setmathfont{Asana-Math.ttf}
% -----------------------------------------------%
% Misc
\author{Κασωτάκη Ε. - Σμαραγδάκης Κ.}
\title{www.math24.gr}
% ---------------- Formation --------------------%


\setlength\topmargin{-0.7cm}
\addtolength\topmargin{-\headheight}
\addtolength\topmargin{-\headsep}
\setlength\textheight{26cm}
\setlength\oddsidemargin{-0.54cm}
\setlength\evensidemargin{-0.54cm}
%\setlength\marginparwidth{1.5in}
\setlength\textwidth{17cm}

\RequirePackage[avantgarde]{quotchap}
\renewcommand\chapterheadstartvskip{\vspace*{0\baselineskip}}
\RequirePackage[calcwidth]{titlesec}
\titleformat{\section}[hang]{\bfseries}
{\Large\thesection}{12pt}{\Large}[{\titlerule[0.9pt]}]
%--------------------------------------------------%
\usepackage{draftwatermark}
\SetWatermarkText{www.math24.gr}
\SetWatermarkLightness{0.9}
\SetWatermarkScale{0.6}
%--------------------------------------------------%


\usepackage{xcolor}
\usepackage{amsthm}
\usepackage{framed}
\usepackage{parskip}

\colorlet{shadecolor}{bluesite!20}

\newtheorem*{orismos}{Ορισμός}

\newenvironment{notation}
  {\begin{shaded}\begin{theorem}}
  {\end{theorem}\end{shaded}}




% Document begins
\begin{document}
\pagestyle{fancy}
\fancyhead{}
\fancyfoot{}
\renewcommand{\headrulewidth}{0pt}
\renewcommand{\footrulewidth}{0pt}


\fancyhead[LO,LE]{
 %\textblockcolor{backsite}
 %\begin{textblock}{5}(-2,-0.55)
  %\rule{0cm}{1cm}
 %\end{textblock}
 \textblockcolor{bluesite}
 \begin{textblock}{5}(-1.5,-0.55)
  \rule{0cm}{1cm}
 \end{textblock}
 %\textblockcolor{bluesite}
 %\begin{textblock}{14}(5.5,-0.55)
 % \rule{0cm}{1cm}
 %\end{textblock}
 \begin{textblock}{0}(-1,-0.25)
 \color{cwhite} \begin{Large}www.math24.gr\end{Large}
 \end{textblock}
\begin{textblock}{0}(16,-0.4)
 \color{cwhite} \includegraphics[height=1.5cm]{math24_logo.png}
 \end{textblock}
\textblockcolor{white}
\begin{textblock}{17}(-1,28)
 \color{backsite} \begin{small}Copyright \textcopyright 2011, Κασωτάκη Ε.(ikasotaki@gmail.com) - Σμαραγδάκης Κ.(kesmarag@gmail.com)\end{small}
 \end{textblock}
}
 \begin{shaded}
\begin{center}
\huge \textbf{Μαθηματικά Α' Γυμνασίου}\\
\end{center} 
\begin{center}
\textbf{"Επαναληπτικές Ασκήσεις για τις διακοπές των Χριστουγέννων"} \hfill \textbf{}
\end{center}
%\subsection*{}
\end{shaded}
%\vspace{2em
%\begin{shaded}
%\begin{center}
%\huge \textbf{Επαναληπτικές Ασκήσεις για το Κεφάλαιο 1}\\
%Πρόσθεση ρητών αριθμών
%\end{center} 
%\textbf{Μαθηματικά Α' Γυμνασίου} \hfill \textbf{Ημερομηνία Παράδοσης : \hspace{2em} }
%\subsection*{Ονοματεπώνυμο :\hfill  \hspace{5em}}
%\end{shaded}
\vspace{2em}
\begin{itemize}
\item Κεφάλαιο 1
\item Κεφάλαιο 7 (7.1-7.5)
\item Κεφάλαιο 2 (2.1-2.2)
\item Κεφάλαιο Β.1 (Β.1.1-Β.1.3)
\end{itemize}


\section*{Άσκηση 1  \hfill \small{}}
Να βρείτε αν οι παρακάτω αριθμοί διαιρούνται με $2,3,4,5$
\begin{enumerate}[i)]
 \item $271$
 \item $444$
 \item $375$
 \item $920$
 \item $111$
\end{enumerate}


\section*{Άσκηση 2  \hfill \small{}}
Να βρείτε το $ΕΚΠ$ και το $ΜΚΔ$  των παρακάτω αριθμών:
\begin{enumerate}[i)]
 \item $15,45$
 \item $8,36$
 \item $13,27$
 \item $8,9$
 \item $18,36$
\end{enumerate}


\section*{Άσκηση 3  \hfill \small{}}
Να αναλυθούν οι αριθμοί $250$ και $426$ σε γινόμενα πρώτων παραγόντων. Με την βοήθεια της ανάλυσης αυτής να βρεθούν 
ο $ΜΚΔ$ και το $ΕΚΠ$ αυτών των αριθμών.


\section*{Άσκηση 4  \hfill \small{}}
Να βρείτε την απόλυτη τιμή των παρακάτω αριθμών: 
\begin{enumerate}[i)]
 \item $-27$
 \item $45$
 \item $0$
 \item $28$
 \item $-108$
\end{enumerate}



\section*{Άσκηση 5  \hfill \small{}}
Να βρείτε τους αριθμούς που έχουν ως απόλυτη τιμή 
\begin{enumerate}[i)]
 \item $78$
 \item $49$
 \item $1072$
 \item $775$
 \item $0$
\end{enumerate}



\section*{Άσκηση 6  \hfill \small{}}
Να γράψετε τους αριθμούς $-3,-8,40,23,-10$ σε αύξουσα σειρά.

\section*{Άσκηση 7  \hfill \small{}}
Να υπολογίσετε την τιμή των παρακάτω αριθμητικών παραστάσεων με δύο τρόπους.\\
(Υπόδειξη: 1ος τρόπος: με χρήση της επιμεριστικής ιδιότητας και 2ος τρόπος: με προτεραιότητα των πράξεων)
\begin{enumerate}[i)]
 \item $-3\cdot (10-2-7)+5\cdot (-1+2+20)$
 \item $(10+5)\cdot (20-10)$
 \item $(3-5)\cdot (10-4)$
 \item $(10+2)\cdot (|-7|-1+2)$
 \item $(4-2)\cdot (2-|-1|+4)$
\end{enumerate}



\section*{Άσκηση 8  \hfill \small{}}
Να κάνετε τις πράξεις χρησιμοποιώντας το αντίστροφο της επιμεριστικής ιδιότητας
\begin{enumerate}[i)]
 \item $3\cdot 17 +3\cdot 3 +3\cdot 71+3\cdot 9$
 \item $2\cdot x+x\cdot 5-4\cdot x -10\cdot x$
 \item $19\cdot 2+19\cdot 6-19\cdot 3+19\cdot 5$
 \item $2\cdot y-13\cdot y -6\cdot y+4\cdot y$
 \item $15\cdot 9-15\cdot 2+3\cdot 15+15\cdot 90$
\end{enumerate}


\section*{Άσκηση 9  \hfill \small{}}
Μία κασετίνα αποτελείται αποκλειστικά από στυλό (μπλε, κόκκινα, μαύρα κ.τ.λ). Τα $\dfrac{3}{8}$ της κασετίνας 
είναι μπλε στυλό. Να βρείτε πόσα στυλό έχει η κασετίνα αν γνωρίζεται ότι τα μπλε στυλό είναι 12.

\section*{Άσκηση 10  \hfill \small{}}
Να εξετάσετε αν τα παρακάτω κλάσματα είναι ισοδύναμα 
\begin{enumerate}[i)]
 \item $\dfrac{11}{22}$ και $\dfrac{2}{4}$
 \item $\dfrac{3}{5}$ και $\dfrac{10}{3}$
 \item $\dfrac{1}{2}$ και $\dfrac{1}{8}$
 \item $\dfrac{3}{5}$ και $\dfrac{9}{15}$
 \item $\dfrac{2}{8}$ και $\dfrac{6}{24}$
\end{enumerate}

\section*{Άσκηση 11  \hfill \small{}}
Να απλοποιήσετε τα παρακάτω κλάσματα. \\
(Υπόδειξη: Να τα μετατρέψετε σε ανάγωγα)
\begin{enumerate}[i)]
 \item $\dfrac{12}{44}$
 \item $\dfrac{3}{9}$
 \item $\dfrac{16}{32}$
 \item $\dfrac{25}{15}$
 \item $\dfrac{55}{35}$
\end{enumerate}



\section*{Άσκηση 12  \hfill \small{}}
Να μετατρέψετε τα παρακάτω κλάσματα σε ομώνυμα
\begin{enumerate}[i)]
 \item $\dfrac{1}{2}$ και $\dfrac{4}{3}$
 \item $\dfrac{2}{3}$ και $\dfrac{10}{4}$
 \item $\dfrac{1}{5}$ και $\dfrac{1}{6}$
 \item $\dfrac{2}{5}$ και $\dfrac{4}{20}$
 \item $\dfrac{3}{2}$ και $\dfrac{2}{4}$
\end{enumerate}


\section*{Άσκηση 13  \hfill \small{}}
Να μετατρέψετε σε $mm$ τα παρακάτω μεγέθη
\begin{enumerate}[i)]
 \item $23cm$
 \item $19dm$
 \item $2m$
 \item $12m$
 \item $9cm$
\end{enumerate}



\end{document}
